\section{Weekly Summary}

\begin{tcolorbox}[colback=yellow!10,colframe=black!50,title=AI-Generated Content]
\noindent The summary below was written by Claude (Anthropic) from that week's regression outputs and series descriptions. It has not been reviewed or fact-checked by a human, and should be read as an automated, informal recap -- not analysis.
\end{tcolorbox}

Welcome back to the little corner of the internet where we take two economic time series that have never met, force them into a spreadsheet together, and see if they hit it off. As always, the house rule stands: nothing here means anything, spurious correlation is the whole point, and if you use any of this to make a decision you deserve the outcome you get. Let's tour the wreckage.

Monday gave us New England commercial electricity producer prices versus a medical price index for, and I quote, "symptoms, signs, and ill-defined conditions." The raw regression came out swinging with an R-squared of 0.93, which is the kind of number that makes you briefly believe you've discovered something. Then we detrended it and the relationship evaporated so completely that the R-squared dropped to a heartwarming 0.02 with a p-value of basically one. In other words, both things went up over the years, and that's it. A textbook trend mirage. The medical category is literally named "ill-defined conditions," and honestly that's the most self-aware series we've ever regressed.

Tuesday paired active home listings in Flathead County, Montana with U.S. imports of maintenance and repair services. The raw regression was weak-ish (R-squared 0.12), which normally means we all go home disappointed. But plot twist: the detrended version got *stronger*, climbing to 0.19 and staying significant. This is the rare, genuinely-worth-a-raised-eyebrow case where the wiggles line up even after the shared drift is stripped out. Do I have a theory? Sure: absolutely none. Montana real estate agents fixing broken things they import? I refuse to explain it, and neither should you.

Wednesday's contestants were wholesale trade payrolls and the Dover, Delaware house price index, and this one actually holds up in both rounds -- raw R-squared of 0.73, detrended R-squared of 0.73, both wildly significant. It's practically the same number twice, which for this project is downright suspicious in a good way. Payrolls and house prices both responding to the general state of the economy is at least a story a grown adult could tell with a straight face, so I'll allow it. Thursday, meanwhile, closed us out on a properly ridiculous note: Southwest-region central-bank GDP versus a discontinued Pacific-division bank-assets ratio. The raw result limped in barely-not-significant (p around 0.055, the statistical equivalent of "we need to talk"), but detrending flipped it into significance at 0.008. Both regressions came with a stern warning about numerical problems and enormous condition numbers, which is the software's polite way of saying "are you sure about this." I am not. I never am. See you next week for more of the same.
